%! Author = sriadityavedantam
%! Date = 3/30/26

\section{Definition of Groups}

A group is a one-operation analogue of a ring.
Groups are often thought of as bijections from a set to itself, or permutations.
Most of the course we will be looking at finite groups.Let $X$ be a set.
Assume $X\neq \varnothing$.
Let $S_X$, called the symmetric group on $X$, be the set of permutations of $X$, equivalently the bijections $X\to X$.Composition has the following properties.
There is an identity permutation, denoted by $e$, $i$, or $1$.
Composition is associative:
\[
    (a\circ b)\circ c = a\circ (b\circ c).
\]
Every element has a unique inverse, since each permutation is a bijection.
If $f$ is a permutation of $X$, then its unique inverse is denoted by $f^{-1}$.An abstract group $G$ is a nonempty set with a binary operation, usually denoted by $\cdot$, $*$, or juxtaposition.
There is an identity element, say $1$, such that
\[
    1*a = a*1 = a
\]
for all $a\in G$.
A group does not need to be commutative, though the identity commutes with every element.
The operation is associative, so for all $a,b,c\in G$,
\[
    (a*b)*c = a*(b*c) = a*b*c.
\]
Every element $a\in G$ has an inverse $a^{-1}$.

\begin{definition}[Abelian Groups]
    A group $G$ does not need to be commutative, but if it is, we call it a commutative group, or \textbf{abelian}.
\end{definition}

$S_X$ is the group of permutations on $X$ under composition.
If $G$ is a subset of $S_X$ which is closed under composition and inverses, then $G$ is a group.
Cayley's Theorem says that every group is isomorphic to a subgroup of some symmetric group.
We will prove this later.Suppose
\[
    X\coloneqq \{1,2,\dots,n\}.
\]
Then $S_X$ is denoted by $S_n$, the symmetric group on $n$ letters.
For example, $S_3$ is the set of permutations of $\{1,2,3\}$.
The notation for permutations is as follows.
Suppose $\sigma\in S_3$.
Then we denote it by
\[
    \begin{pmatrix}
        1&2&3\\
        \sigma(1)&\sigma(2)&\sigma(3)
    \end{pmatrix}.
\]
The elements of $S_3$ are
\begin{gather*}
    e=
    \begin{pmatrix}
        1&2&3\\
        1&2&3
    \end{pmatrix},
    \quad
    \sigma_1=
    \begin{pmatrix}
        1&2&3\\
        1&3&2
    \end{pmatrix},
    \quad
    \sigma_2=
    \begin{pmatrix}
        1&2&3\\
        2&1&3
    \end{pmatrix},\\
    \sigma_3=
    \begin{pmatrix}
        1&2&3\\
        2&3&1
    \end{pmatrix},
    \quad
    \sigma_4=
    \begin{pmatrix}
        1&2&3\\
        3&2&1
    \end{pmatrix},
    \quad
    \sigma_5=
    \begin{pmatrix}
        1&2&3\\
        3&1&2
    \end{pmatrix}.\\
\end{gather*}
Suppose $\sigma_2\circ \sigma_4$.
Then we follow elements through $\sigma_4$ first and then $\sigma_2$.
Doing this gives
\[
    \begin{pmatrix}
        1&2&3\\
        3&1&2
    \end{pmatrix}.
\]
\begin{example}
    \[
        \sigma_4\circ \sigma_2 =
        \begin{pmatrix}
            1&2&3\\
            2&3&1
        \end{pmatrix}.
    \]
\end{example}

If $(\mathcal{R},+,\cdot)$ is a ring with operations addition and multiplication, then $(\mathcal{R},+)$ is a group.
The identity in an additive group is $0$.
It is associative, and every element $a$ has an additive inverse, denoted by $-a$.
In fact, these groups are abelian.
Let $\mathbb{F}$ be any field.
Let
\[
    \mathbb{F}^* \coloneqq \{a\in \ff : a\neq 0\} \coloneqq U(\mathbb{F}).
\]
Then $\mathbb{F}^*$ is a group under multiplication.
This group has identity $1\in \mathbb{F}$, it is associative, and every element has an inverse.
More generally, if $\mathcal{R}$ is a ring, then
\[
    \mathcal{R}^* \coloneqq \{a\in \mathcal{R} : \text{$a$ is a unit in }\mathcal{R}\} \coloneqq U(\mathcal{R})
\]
is a group.

\begin{example}
    \[
        \mathbb{Z}_6^* = \{1,5\}.
    \]
\end{example}

